\documentclass[11pt,a4paper]{article}
\usepackage[margin=20mm]{geometry}
\usepackage{amsmath}
\usepackage[hidelinks]{hyperref}
\setlength{\parindent}{0pt}
\setlength{\parskip}{6pt}
\raggedright\emergencystretch=3em
\begin{document}
\begin{center}
{\Large Why the donor-capacity options give different amounts}\\[3pt]
9 October 2026 --- a short explanation of the existing gate example
\end{center}

\textbf{The setting.} Consider one 1-second update. Traffic can leave a donor
lane either by continuing through or by changing lane. The through destination
is blocked; the lane-change destination has enough receiving capacity. Traffic
reaches the donor's exit only after 0.9 seconds, leaving \textbf{0.1 seconds}
to discharge. All amounts are continuous vehicle-equivalents (eq), not integer
vehicles.

\textbf{Where the shared numbers come from.} The 30-m donor initially contains
traffic in its first 3 m, at density 0.15 eq/m. The leading front covers the
remaining 27 m at 30 m/s. The example's physical sending capacity is
$Q=9/14\simeq0.642857$ eq/s. There is no earlier exit flow.
\[
\begin{aligned}
\text{initial traffic}&=3\times0.15=0.45\ \mathrm{eq},\\
\text{arrival time}&=(30-3)/30=0.9\ \mathrm{s},\\
\text{lane-change allowance}&=25\%\times0.45=0.1125\ \mathrm{eq}.
\end{aligned}
\]
The 25\% is the original sampled lane-choice fraction. The 0.1125-eq allowance
limits \emph{total lane-changing mass during this update}. Both options keep it.

\textbf{D1: reserve an instantaneous capacity fraction for each movement.}
The 25\% fraction also caps the lane-change \emph{rate}. Blocking the through
exit does not release its fraction for lane changing. Thus:
\[
\begin{aligned}
\text{lane-change rate}&=0.25\times0.642857\simeq0.160714\ \mathrm{eq/s},\\
\text{accepted mass}&=\min\{0.160714\times0.1,\ 0.1125\}
\simeq\boxed{0.016071\ \mathrm{eq}}.
\end{aligned}
\]
\textbf{Core idea:} the original direction fraction governs both how much
traffic may change lane and how quickly it may be served. Some donor sending
capability stays unused when the through exit is blocked.

\textbf{D2: share physical sending among still-active original requests.}
Lane changing is the only usable exit, so it can use the donor's full physically
available sending, until its original amount allowance is exhausted:
\[
\text{accepted mass}=\min\{0.642857\times0.1,\ 0.1125\}
\simeq\boxed{0.064286\ \mathrm{eq}}.
\]
\textbf{Core idea:} the fraction limits total requested mass, without imposing
an additional direction-specific rate cap. This borrows unused \emph{sending
capability}, not the blocked movement's mass allowance. It creates no new
lane-changing demand.

\textbf{Why the difference matters.} Both results are below 0.1125 eq, so both
respect the same amount limit. If the existing accepted-entry counter is 0.94,
the ending counters are:
\[
\text{D1: }0.94+0.016071\simeq0.956071<1,
\qquad
\text{D2: }0.94+0.064286\simeq1.004286>1.
\]
Under the already approved E1 bookkeeping rule, only D2 therefore generates a
new restriction record at the next regular update. That record adds no traffic
mass. This shows a real difference in the forward laws; it does not prove that
D2 is more realistic or gives a better control policy. Neither option is adopted.

\small
For the exact underlying traffic fields, see Section 6 of
\href{fixed_share_completion_math.pdf}{the existing joint-example derivation}.
Here rates are rounded for readability; the exact accepted amounts are
$9/560$ eq (D1) and $9/140$ eq (D2).
\end{document}
